\section*{الفصل \ref{ch:b1:elimination} --- الحذف \texorpdfstring{$\beta$}{بيتا}}

\begin{solution}{exo:b1:elimination:1}
2-بروموبوتان: بوت-1-ين، وبوت-2-ين-\cip{E} و-\cip{Z}؛ والرئيسي
بوت-2-ين-\cip{E}. 2-برومو-2-ميثيل البوتان: 2-ميثيل بوت-2-ين (الرئيسي،
ثلاثي الاستبدال) و2-ميثيل بوت-1-ين. البروموهكسان الحلقي: الهكسين الحلقي وحده.
\end{solution}

\begin{solution}{exo:b1:elimination:2}
E1: $v = k[\mathrm{RBr}]$؛ E2: $v = k[\mathrm{RBr}][\ce{EtO-}]$. قِس
\omterm{def:b1:rate-laws:initial-rate}{السرعة الابتدائية} لتكوّن الألكين عند عدة تراكيز للإيثوكسيد: فهي
تزداد بالتناسب في حالة E2، ولا تتغير في حالة E1.
\end{solution}

\begin{solution}{exo:b1:elimination:3}
على امتداد C2--C3، يمكن أن تكون ذرة Br على C2 مضادة لأيٍّ من هيدروجيني C3 أو
لمجموعة الميثيل C4: والتشكّلان اللذان تكون فيهما H مضادة لذرة Br يمكن أن
يتفاعلا \omterm{def:b1:elimination:elimination}{بالآلية E2} (فيعطيان الألكينين \cip{E} و\cip{Z})؛ أما الثالث فلا يمكنه
الحذف نحو C3.
\end{solution}

\begin{solution}{exo:b1:elimination:4}
SN2 (أولي، هيدروكسيد، درجة حرارة منخفضة)؛ E2 (ثالثي، \omterm{def:b1:predominant-reaction:strong-acid}{قاعدة قوية}
ضخمة)؛ SN2 (اليوديد، \omterm{def:b1:electronic-effects:nucleophile}{محب للنواة} جيد \omterm{def:b1:predominant-reaction:strong-acid}{وقاعدة ضعيفة}، \omterm{def:b1:intermolecular-forces-solvents:solvent-classes}{مذيب لابروتوني})؛
SN1 وE1، وتزداد E1 بالتسخين.
\end{solution}

\begin{solution}{exo:b1:elimination:5}
يحمل C3 هيدروجينًا واحدًا. وفي كل \omterm{def:b1:stereochemistry-in-depth:enantiomers}{متماكب لامرآوي}، يحدد وضعُ ذلك الهيدروجين مضادًا
لذرة Br على C2 مواضعَ المجموعات الأخرى: فميثيل C2
وإيثيل C3 في الجهة نفسها في أحد \omterm{def:b1:stereochemistry-in-depth:enantiomers}{المتماكبين اللامرآويين} وفي
جهتين متقابلتين في الآخر. فيعطي أحدهما 3-ميثيل بنت-2-ين-\cip{E}، والآخر
3-ميثيل بنت-2-ين-\cip{Z}: وهذا \omterm{def:b1:nucleophilic-substitution:stereospecific}{تفاعل نوعي فراغيًا}.
\end{solution}

\begin{solution}{exo:b1:elimination:6}
تتبرتن مجموعة OH، \ce{R-OH2+}؛ ويغادر الماء، فيتكوّن \omterm{def:b1:electronic-effects:intermediates}{الكاتيون الكربوني}
الثالثي \ce{(CH3)2C+-CH2CH3}؛ وتنزع قاعدة (الماء، \ce{HSO4-})
بروتونًا من كربون مجاور. الناتج الرئيسي (زايتسيف):
2-ميثيل بوت-2-ين.
\end{solution}

\begin{solution}{exo:b1:elimination:7}
تثبّت مجموعة ثالثي البوتيل \omterm{def:b1:stereochemistry-in-depth:conformations}{الكرسي} الذي تكون فيه هي \omterm{def:b1:stereochemistry-in-depth:conformations}{استوائية}. وفي
المتماكب \textit{trans} يكون البروم عندئذ استوائيًا أيضًا، فلا يكون أبدًا مضادًا لرابطة
C--H: فعلى E2 أن تنتظر \omterm{def:b1:stereochemistry-in-depth:conformations}{الكرسي} المقلوب غير الملائم جدًا. وفي
المتماكب \textit{cis} يكون البروم محوريًا، ومعه هيدروجينان ثنائيا المحور
\textit{trans}: فتكون E2 سريعة.
\end{solution}

\begin{solution}{exo:b1:elimination:8}
الإيثوكسيد، الصغير، ينزع الهيدروجين الذي يعطي الألكين الأكثر
استبدالًا (زايتسيف). أما ثالثي البوتوكسيد، الضخم، فيبلغ الهيدروجينات المكشوفة
لمجموعة الميثيل بسهولة أكبر من هيدروجين مجموعة \ce{CH2}
الداخلية: فيعطي الألكين الأقل استبدالًا.
\end{solution}

\begin{solution}{exo:b1:elimination:9}
ليس للبروموميثان كربون $\beta$. وفي 1-برومو-2,2-ثنائي ميثيل البروبان يكون
الكربون $\beta$ رباعيًا ولا يحمل أي هيدروجين.
\end{solution}

\begin{solution}{exo:b1:elimination:10}
يحتاج الناتجان كلاهما إلى \omterm{def:b1:electronic-effects:intermediates}{الكاتيون الكربوني}، المتكوّن في الخطوة البطيئة بالسرعة
$k_1[\mathrm{RBr}]$. ثم يتوزع بين الأسر بالإيثانول (\omterm{def:b1:rate-laws:order}{ثابت
السرعة} $k_S$) وفقد بروتون (\omterm{def:b1:rate-laws:order}{ثابت السرعة} $k_E$)، وكلاهما
من رتبة أولى كاذبة بالنسبة إلى المذيب: فكسر الألكين هو
$k_E/(k_E + k_S)$، مستقل عن تركيز \omterm{def:b1:nucleophilic-substitution:substitution}{الركيزة}.
\end{solution}

\begin{solution}{exo:b1:elimination:11}
انطلاقًا من \omterm{def:b1:stereochemistry-in-depth:isomers}{تشكّل} \omterm{def:b1:stereochemistry-in-depth:enantiomers}{مركب الميزو} الذي تكون فيه ذرتا Br متضادتين
(فتكون مجموعتا الفينيل متضادتين، وذرتا H متضادتين)، دوّر C1 بزاوية $120^\circ$ لتضع
هيدروجينه مضادًا لذرة Br على C2. فتصير مجموعتا الفينيل عندئذ في
الجهة نفسها من المحور H--Br: فتنتهيان في وضع \textit{cis}. وعلى C1، يتقدم Br على
Ph؛ وعلى C2، يتقدم Ph على H؛ وBr وفينيل C2 في وضع \textit{trans}:
1-برومو-1,2-ثنائي فينيل الإيثين-\cip{E}.
\end{solution}

\begin{solution}{exo:b1:elimination:12}
لا يتفاعل إلا \omterm{def:b1:stereochemistry-in-depth:conformations}{الكرسي} الذي تكون فيه \omterm{def:b1:electronic-effects:nucleophile}{المجموعة المغادرة} \omterm{def:b1:stereochemistry-in-depth:conformations}{محورية}: $v = k\,x\,[\mathrm{RY}][\mathrm{B}]$،
حيث $x$ كسر ذلك \omterm{def:b1:stereochemistry-in-depth:conformations}{الكرسي}. وكل مجموعة ألكيل \omterm{def:b1:stereochemistry-in-depth:conformations}{محورية} يتطلبها
تكلّف عدة kJ/mol (\qty{5.7}{kJ/mol} لمجموعة الميثيل)، فتقسم $x$ على نحو
عشرة لكل مجموعة عند درجة حرارة الغرفة: ومع مجموعتين أو ثلاث من هذا النوع، تهبط $x$
إلى $10^{-2}$ أو $10^{-3}$ ويصير التفاعل أبطأ بالقدر نفسه.
\end{solution}

\begin{solution}{pb:b1:elimination:1}
\textbf{1.} \omterm{def:b1:stereochemistry-in-depth:conformations}{كرسي} تكون فيه Cl والإيزوبروبيل والميثيل كلها \omterm{def:b1:stereochemistry-in-depth:conformations}{استوائية}.
\textbf{2.} الإيزوبروبيل والميثيل استوائيان، وCl \omterm{def:b1:stereochemistry-in-depth:conformations}{محورية}.
\textbf{3.} \omterm{def:b1:stereochemistry-in-depth:enantiomers}{متماكبان لامرآويان}: لا يختلفان إلا عند C1.
\textbf{4.} \omterm{def:b1:electronic-effects:nucleophile}{المجموعة المغادرة} وهيدروجين $\beta$ ثنائيا المحور \textit{trans}.
\textbf{5.} في \omterm{def:b1:stereochemistry-in-depth:conformations}{الكرسي} المقلوب وحده، حيث Cl والإيزوبروبيل والميثيل كلها
\omterm{def:b1:stereochemistry-in-depth:conformations}{محورية}.
\textbf{6.} المجموعات الثلاث كلها \omterm{def:b1:stereochemistry-in-depth:conformations}{محورية}: الميثيل وحده يكلّف نحو
\qty{5.7}{kJ/mol}، ومجموعة الإيزوبروبيل الأكبر أكثر، والكلور قليلًا:
فالمجموع نحو \qty{13}{kJ/mol}، ومن ثم كسر يقارب
$\exp(-13000/2479) = \num{5e-3}$، وهي القيمة التي تأخذها المسألة.
\textbf{7.} C2 (هيدروجين واحد، إذ يحمل مجموعة الإيزوبروبيل) وC6 (هيدروجينان).
\textbf{8.} مع Cl \omterm{def:b1:stereochemistry-in-depth:conformations}{محورية} على C1، الهيدروجين \omterm{def:b1:stereochemistry-in-depth:conformations}{المحوري} على C2 (إذ الإيزوبروبيل
\omterm{def:b1:stereochemistry-in-depth:conformations}{استوائي}) والهيدروجين \omterm{def:b1:stereochemistry-in-depth:conformations}{المحوري} على C6: اثنان.
\textbf{9.} في \omterm{def:b1:stereochemistry-in-depth:conformations}{الكرسي} الثلاثي المحور تشغل مجموعة الإيزوبروبيل الموضع \omterm{def:b1:stereochemistry-in-depth:conformations}{المحوري}
على C2، فيكون هيدروجينه استوائيًا: غير مضاد. ولذرة C6 هيدروجين \omterm{def:b1:stereochemistry-in-depth:conformations}{محوري}: مضاد.
\textbf{10.} كلوريد النيومنثيل اثنان، وكلوريد المنثيل واحد.
\textbf{11.} على امتداد C1--C6: Cl (\omterm{def:b1:stereochemistry-in-depth:conformations}{محورية}، إلى الأعلى) على C1 مقابل الهيدروجين \omterm{def:b1:stereochemistry-in-depth:conformations}{المحوري} (إلى الأسفل)
على C6؛ وروابط الحلقة والهيدروجين \omterm{def:b1:stereochemistry-in-depth:conformations}{الاستوائي} متخالفة بينهما.
\textbf{12.} على \omterm{def:b1:stereochemistry-in-depth:conformations}{الكرسي} لا توجد رابطة C--H \omterm{def:b1:elimination:periplanar}{في وضع متزامن مستوٍ} بالنسبة إلى الرابطة C--Cl؛
وفرض ذلك يعني حلقة متطابقة مجهَدة.
\textbf{13.} بنزع هيدروجين C2: 4-ميثيل-1-(بروبان-2-يل)هكس-1-ين حلقي
(ثلاثي الاستبدال)؛ وبنزع هيدروجين C6: 3-ميثيل-6-(بروبان-2-يل)هكس-1-ين حلقي
(ثنائي الاستبدال).
\textbf{14.} الألكين الثلاثي الاستبدال، الأكثر استبدالًا وثباتًا.
\textbf{15.} الألكينان كلاهما، \qty{78}{\percent} منه ثلاثي الاستبدال: متسق.
\textbf{16.} الألكين الثنائي الاستبدال وحده: عكس ما تتنبأ به
\omterm{prop:b1:elimination:zaitsev}{قاعدة زايتسيف}، وقد فرضه الهيدروجين المضاد الوحيد المتاح.
\textbf{17.} كلاهما نوعي فراغيًا (شرط الوضع المضاد)؛ ويتفاعل كلوريد النيومنثيل
تفاعلًا انتقائيًا موضعيًا (78/22)؛ ويعطي كلوريد المنثيل متماكبًا
بنيويًا وحيدًا.
\textbf{18.} الكربون الحامل لذرة Cl ثانوي ومحاط بمجموعة
إيزوبروبيل: مزدحم بالنسبة إلى SN2، بينما الإيثوكسيد \omterm{def:b1:predominant-reaction:strong-acid}{قاعدة قوية}.
\textbf{19.} $v \propto x\,n_{\mathrm{H}}$: كسر \omterm{def:b1:stereochemistry-in-depth:conformations}{الكرسي} النشط مضروبًا في
عدد الهيدروجينات المضادة.
\textbf{20.} $0.95 \times 2 = 1.90$.
\textbf{21.} $\num{5.0e-3} \times 1 = \num{5.0e-3}$.
\textbf{22.} \textbf{$k(\text{نيومنثيل})/k(\text{منثيل}) = 1.90/\num{5.0e-3}
= 380$}.
\end{solution}
